% =============================================================
% 01_strategy.tex  第1部 ストラテジ系
% =============================================================
\chapter{ストラテジ系：何のために作るかを考える}

\section{企業活動と経営戦略}

\pitfall{「経営戦略」と「経営方針」と「事業戦略」の違いがわからない}{
初学者がまずつまずくのが、戦略関連の用語のレイヤーです。
大きい順に、\termtag{経営理念} → \termtag{経営方針} → \termtag{経営戦略} → \termtag{事業戦略} → \termtag{機能戦略} と、
\textbf{抽象から具体へ階段状}に並んでいると押さえてください。
}

\definitionbx{経営戦略}{
企業が\textbf{中長期的に}競争優位を築くために、リソース（人・モノ・金・情報）をどこに集中させるかを決める方針。
同じ業界内でも、差別化・コストリーダーシップ・集中などの取り方がある（\termtag{ポーターの基本戦略}）。
}

\comparisonbx{ビジネスでよく出る分析フレーム}{%
\comptable{フレーム & 対象 & 一言で}{
SWOT & 自社 & 強み/弱み/機会/脅威を整理 \\
3C & 市場 & Customer / Competitor / Company \\
4P & マーケ & Product / Price / Place / Promotion \\
PPM & 事業 & 市場成長率 × 市場占有率で投資判断
}
}

\advice{%
これらは\textbf{「何を見るための道具か」}で区別します。
「自社の健康診断」＝SWOT、「市場の三角測量」＝3C、「販売設計」＝4P、「投資ポートフォリオ」＝PPM と覚えると記憶に残ります。
}

\section{会計・財務の最低限}

\pitfall{「売上＝利益」ではない}{
売上は入ってきたお金、利益はそこから費用を引いた残り。
さらに、売上総利益・営業利益・経常利益・純利益と\textbf{段階ごとに意味が違います}。
損益計算書（PL）は「\termtag{ふるいが下がっていくイメージ}」で読むのがコツです。
}

\comparisonbx{財務三表の役割}{%
\comptable{名前 & 何を表すか & たとえ}{
貸借対照表 (BS) & ある時点の\textbf{残高} & 家計の残高スナップショット \\
損益計算書 (PL) & 期間中の\textbf{儲け} & 1 年分の家計簿 \\
キャッシュフロー計算書 (CF) & 期間中の\textbf{現金の流れ} & 実際の出入金の記録
}
}

\section{法務・知的財産・個人情報}

\pitfall{「個人情報」と「プライバシー」と「機密情報」の違いが曖昧}{
試験では、用語が似ていて紛らわしい三兄弟が何度も登場します。
\begin{itemize}
  \item \termtag{個人情報}：個人を識別できる情報（氏名・住所・顔写真など）
  \item \termtag{プライバシー}：本人が他人に知られたくない領域（法律上の概念としては「個人情報」より広い）
  \item \termtag{機密情報}：組織の秘密として扱う情報（個人に限らない）
\end{itemize}
\textbf{対象と視点が違う}ので、「法律の対象」「個人の権利」「会社の秘密」と言い換えて覚えましょう。
}

\comparisonbx{知的財産権の整理}{%
\comptable{権利 & 保護対象 & 代表的なもの}{
著作権 & 表現 & 文章・音楽・コード \\
特許権 & 発明 & 高度な技術的アイデア \\
実用新案 & 考案 & 小さな工夫・形状 \\
意匠権 & デザイン & 工業製品の形 \\
商標権 & 識別標識 & ロゴ・商品名
}
}

\checkquestion{SWOT 分析で扱うのはどれか。}
  {市場の成長率と占有率}
  {自社の強み・弱みと、外部の機会・脅威}
  {製品・価格・流通・販促の 4 要素}
  {顧客・競合・自社の 3 視点}
  {B}
  {A は PPM、C は 4P、D は 3C の話です。SWOT は「自社の健康診断」と覚えるのが最短です。}

\section{経営戦略のキーワード整理}

\pitfall{横文字フレームワークが多すぎて頭に入らない}{
ストラテジ系は横文字が多く、3 文字や 4 文字の略語が次々に出てきます。
\textbf{「何を見るための道具か」}を決めて引き出しに収めれば、それほど怖くありません。
ここでは頻出フレームワークを「自社を見る／市場を見る／投資判断する」の 3 つの引き出しに整理します。
}

\comparisonbx{戦略・分析フレームワークの引き出し}{%
\comptable{フレーム & 観点 & 一言で}{
SWOT & 自社 & 強み/弱み $\times$ 機会/脅威 \\
VRIO & 自社 & 強みが\textbf{本当に競争優位か}を問う \\
3C & 市場 & 顧客・競合・自社 \\
4P / 4C & マーケ & 売り手 4 視点 / 買い手 4 視点 \\
PPM & 事業 & 成長率 $\times$ 占有率で投資配分 \\
アンゾフ & 成長 & 既存／新規 $\times$ 市場／製品
}
}

\pitfall{「イノベーション」と「DX」と「リーンスタートアップ」の関係}{
\begin{itemize}
  \item \termtag{イノベーション}：新しい価値や仕組みを生むこと（最も広い概念）
  \item \termtag{DX（デジタルトランスフォーメーション）}：IT を使って\textbf{経営やビジネスモデルを変える}こと
  \item \termtag{リーンスタートアップ}：\textbf{最小限の製品}で仮説検証を繰り返す\textbf{手法}
\end{itemize}
階段で言えば、「目的＝イノベーション」「文脈＝DX」「手法＝リーンスタートアップ」の関係です。
}

\comparisonbx{新事業にまつわる用語の棲み分け}{%
\comptable{用語 & 要点 & 似ている用語との違い}{
MVP & 検証用の最小製品 & 本番製品ではない \\
PoC & 技術の実現性確認 & 商業化の前段階 \\
PoV & 価値があることの実証 & PoC より一歩先 \\
キャズム & 一般市場への跳躍時の壁 & 死の谷より後ろ \\
死の谷 & 研究→事業化の資金・人の壁 & 魔の川より後ろ \\
M\&A & 合併・買収 & アライアンス（提携）より深い
}
}

\onelinesummary{%
ストラテジ系は\textbf{「何のために作るか」}を決める分野。
戦略・財務・法務はどれも「\termtag{見る対象}」と「\termtag{視点}」を分けて整理すると混乱しない。
フレームワークは\textbf{用途別の引き出し}に仕分けてから覚える。
}
